\FloatBarrier
\section{Fitted parameters}\label{si:params}

\begin{table}[!htb]
\centering\small
\caption{Fixed-window fits (one-step least squares on the training window). $\phi$ is the lag-one coefficient of the training residuals, clipped to $[-0.95,0.95]$; $\sigma$ their standard deviation (kt). M2--M4 share the stock-flow fit; M1 uses the training-mean catch, M1b the depensatory branch. Bounds: $r\in[10^{-3},2]$, $K\in[\max S+10,\,5000]$, $\sA\in[0,\max S]$; initialiser $r_0=0.3$, $K_0=\max(1.5\max S, 500)$.}
\label{tab:si_fixedfits}
\begin{tabular}{@{}llrrrrr@{}}
\toprule
Window (catch) & Module & $r$ & $K$ (kt) & $\sA$ (kt) & $\phi$ & $\sigma$ \\
\midrule
A collapse, train 1983--90 (coarse) & M1 & 1.935 & 1032.7 & -- & $-0.14$ & 33.8 \\
 & M1b & 2.000 & 1058.9 & 0.0 & $-0.65$ & 40.2 \\
 & M2--M4 & 1.935 & 1032.7 & -- & $-0.14$ & 33.8 \\
A recovery, train 1995--2007 (coarse) & M1 & 0.458 & 500.0 & -- & $0.95$ & 11.8 \\
 & M1b & 2.000 & 105.9 & 0.0 & $0.95$ & 11.7 \\
 & M2--M4 & 0.458 & 500.0 & -- & $0.95$ & 11.8 \\
A collapse (annual landings) & M1 & 1.942 & 1034.3 & -- & $-0.14$ & 33.8 \\
 & M1b & 2.000 & 1062.0 & 0.0 & $-0.66$ & 40.5 \\
 & M2--M4 & 1.991 & 1032.3 & -- & $-0.02$ & 47.0 \\
A recovery (annual landings) & M1 & 0.370 & 5000.0 & -- & $0.95$ & 11.8 \\
 & M1b & 2.000 & 129.8 & 0.0 & $0.95$ & 11.6 \\
 & M2--M4 & 0.390 & 5000.0 & -- & $0.95$ & 10.6 \\
\bottomrule
\end{tabular}
\end{table}

The $K$ estimate of the coarse-regime recovery fit, 500.0, is the initialiser $K_0 = \max(1.5\max S, 500)$: the objective is flat in $K$ above the observed range and the single L-BFGS-B start does not move it. Under annual landings the same fit runs to the 5000-kt upper bound. In both cases the production curve over the observed range (10--81\,kt) is effectively linear with slope $r$ (main-text Fig.~4).

\begin{table}[!htb]
\centering\scriptsize
\setlength{\tabcolsep}{4pt}
\caption{Specification A, coarse regime: parameters refitted at every rolling origin (expanding window from 1983). $\bar C$ is the training-mean catch used by M1 and M1b. From the 1992 origin, once the collapse transitions enter the training window, $K$ sits at its bound for the stock-flow fit and the depensatory threshold moves to 585--734\,kt; from 2011 it falls to about 5\,kt.}
\label{tab:si_rollingA}
\begin{tabular}{@{}lrrrrrrrrr@{}}
\toprule
 & \multicolumn{3}{c}{M2--M4 (stock-flow)} & \multicolumn{3}{c}{M1 (autonomous)} & \multicolumn{3}{c}{M1b (depensatory)} \\
\cmidrule(lr){2-4}\cmidrule(lr){5-7}\cmidrule(lr){8-10}
Origin & $r$ & $K$ & $\phi$ & $\bar C$ & $r$ & $K$ & $r$ & $K$ & $\sA$ \\
\midrule
1990 & 1.935 & 1033 & -0.14 & 240.0 & 1.935 & 1033 & 2.00 & 1059 & 0.0 \\
1991 & 1.557 & 1058 & -0.03 & 240.0 & 1.557 & 1058 & 2.00 & 1042 & 0.0 \\
1992 & 0.269 & 5000 & +0.34 & 228.0 & 0.252 & 5000 & 2.00 & 1192 & 672.7 \\
1993 & 0.250 & 5000 & +0.67 & 207.7 & 0.212 & 5000 & 2.00 & 2110 & 635.0 \\
1994 & 0.248 & 5000 & +0.57 & 190.8 & 0.190 & 5000 & 2.00 & 3049 & 585.3 \\
1995 & 0.248 & 5000 & +0.56 & 176.5 & 0.170 & 5000 & 2.00 & 3300 & 597.9 \\
1996 & 0.248 & 5000 & +0.56 & 164.3 & 0.153 & 5000 & 2.00 & 3271 & 628.8 \\
1997 & 0.248 & 5000 & +0.56 & 153.7 & 0.138 & 5000 & 2.00 & 3332 & 649.2 \\
1998 & 0.248 & 5000 & +0.56 & 144.4 & 0.125 & 5000 & 2.00 & 3453 & 663.6 \\
1999 & 0.248 & 5000 & +0.56 & 136.2 & 0.114 & 5000 & 2.00 & 3633 & 673.3 \\
2000 & 0.248 & 5000 & +0.56 & 128.9 & 0.104 & 5000 & 2.00 & 3942 & 676.1 \\
2001 & 0.248 & 5000 & +0.56 & 122.4 & 0.095 & 5000 & 2.00 & 4198 & 681.7 \\
2002 & 0.304 & 2537 & +0.58 & 116.5 & 0.087 & 5000 & 2.00 & 4373 & 690.6 \\
2003 & 0.247 & 5000 & +0.56 & 111.2 & 0.080 & 5000 & 2.00 & 4476 & 701.6 \\
2004 & 0.247 & 5000 & +0.56 & 106.4 & 0.073 & 5000 & 2.00 & 4547 & 713.0 \\
2005 & 0.247 & 5000 & +0.56 & 102.0 & 0.067 & 5000 & 2.00 & 4606 & 723.7 \\
2006 & 0.248 & 5000 & +0.56 & 97.9 & 0.061 & 5000 & 2.00 & 4768 & 730.7 \\
2007 & 0.248 & 5000 & +0.56 & 94.2 & 0.057 & 5000 & 2.00 & 5000 & 734.1 \\
2008 & 0.248 & 5000 & +0.56 & 90.8 & 0.054 & 5000 & 1.57 & 5000 & 707.8 \\
2009 & 0.248 & 5000 & +0.56 & 87.6 & 0.050 & 5000 & 1.12 & 5000 & 656.9 \\
2010 & 0.247 & 5000 & +0.56 & 84.6 & 0.047 & 5000 & 0.80 & 5000 & 588.7 \\
2011 & 0.246 & 5000 & +0.56 & 81.9 & 0.044 & 5000 & 0.13 & 1411 & 5.3 \\
2012 & 0.247 & 5000 & +0.56 & 79.3 & 0.042 & 5000 & 0.12 & 1411 & 5.3 \\
2013 & 0.247 & 5000 & +0.56 & 76.9 & 0.083 & 1408 & 0.11 & 1411 & 5.2 \\
2014 & 0.249 & 5000 & +0.56 & 74.7 & 0.148 & 1114 & 0.11 & 1411 & 5.1 \\
\bottomrule
\end{tabular}
\end{table}

\begin{table}[!htb]
\centering\scriptsize
\setlength{\tabcolsep}{4pt}
\caption{Specification B: range of the refitted parameters over the origins of each band. The depensatory threshold stays at or near zero throughout; the stock-flow $K$ is at or near its bound after 1993.}
\label{tab:si_rollingB}
\scriptsize\setlength{\tabcolsep}{2.5pt}
\begin{tabular}{@{}lrrrrrrrrr@{}}
\toprule
 & & \multicolumn{2}{c}{M2--M4} & \multicolumn{3}{c}{M1} & \multicolumn{3}{c}{M1b} \\
\cmidrule(lr){3-4}\cmidrule(lr){5-7}\cmidrule(lr){8-10}
Origins & $n$ & $r$ & $K$ & $\bar C$ & $r$ & $K$ & $r$ & $K$ & $\sA$ \\
\midrule
1965--1990 & 26 & 0.46--0.78 & 2482--5000 & 354--468 & 0.77--0.95 & 1845--2276 & 1.41--2.00 & 1642--2279 & 0.00--0.27 \\
1991--1995 & 5 & 0.50--0.59 & 3355--4800 & 317--349 & 0.71--0.83 & 1869--1922 & 1.57--1.76 & 1654--1660 & 0.00--0.00 \\
1996--2012 & 17 & 0.50--0.50 & 4802--5000 & 227--310 & 0.56--0.70 & 1789--1905 & 0.91--1.53 & 1651--2261 & 0.00--0.04 \\
2013--2023 & 11 & 0.47--0.50 & 5000--5000 & 193--223 & 0.53--0.57 & 1682--1774 & 0.81--1.09 & 1603--2261 & 0.00--0.04 \\
\bottomrule
\end{tabular}
\end{table}

\begin{remark}[The compensatory branch is not the $\sA\to 0$ limit of the depensatory one]
Substituting $\sA = 0$ in $a(S) = (S-\sA)/(K-\sA)$ gives $a(S) = S/K$, so $g(S) = rS(1-S/K)(S/K)$, a cubic surplus, not the Schaefer form $a\equiv 1$. M1b at $\sA = 0$ is therefore a different model from M1, which is why the two differ slightly in the collapse window (693.6 against 694.3\,kt) even though both thresholds are at zero; what M1b adds to M1 is a branch of the family, not a limiting case of it.
\end{remark}

\begin{figure}[!htb]
\centering
\includegraphics[width=\textwidth]{figS1_allee_profile}
\caption{Profile of the one-step objective over a fixed depensation threshold $\sA$ in the recovery window (train 1995--2007), with $r$ and $K$ re-optimised at each value. (a) The objective rises monotonically from $\sA = 0$ and changes by 6.7\% (coarse regime, 126.5 to 135.0\,kt$^2$) and 2.9\% (annual landings, 122.9 to 126.5) across the range below the smallest observed state (9.7\,kt). (b) The compensating $K$ falls as $\sA$ rises, with $r$ at its bound of 2 throughout. The threshold is unidentified in this window; its least-squares value is the boundary.}
\label{fig:si_profile}
\end{figure}

\FloatBarrier
\section{Floor-forecast census}\label{si:floor}

A forecast is at the floor when the iterated map reaches the numerical floor $\epsilon = 10^{-3}\kt$; its absolute error then equals the observed SSB. Table~\ref{tab:si_floor} counts floor forecasts by band from the archived per-origin files (coarse regime for Specification A). The counts in the first block are those behind the superscripts of Table 6 and Table~\ref{tab:si_h5}.

\begin{table}[!htb]
\centering\small
\caption{Number of origins at which the forecast is at the floor, by band and horizon.}
\label{tab:si_floor}
\begin{tabular}{@{}llrrrrrr@{}}
\toprule
Spec. & Origins & $n$ & M1 & M1b & M2 & M3 & M4 \\
\midrule
A, $h=1$ & 1990--1995 & 6 & 3 & 4 & 0 & 3 & 0 \\
A, $h=1$ & 1996--2012 & 17 & 12 & 13 & 0 & 0 & 0 \\
A, $h=1$ & 2013--2014 & 2 & 0 & 0 & 0 & 0 & 0 \\
A, $h=5$ & 1990--1995 & 6 & 4 & 4 & 1 & 3 & 1 \\
A, $h=5$ & 1996--2010 & 15 & 15 & 15 & 0 & 0 & 0 \\
B, $h=1$ & 1965--1990 & 26 & 2 & 3 & 0 & 0 & 0 \\
B, $h=1$ & 1991--1995 & 5 & 3 & 3 & 0 & 3 & 1 \\
B, $h=1$ & 1996--2012 & 17 & 17 & 17 & 0 & 0 & 0 \\
B, $h=1$ & 2013--2023 & 11 & 0 & 1 & 0 & 0 & 0 \\
B, $h=5$ & 1965--1990 & 26 & 10 & 17 & 7 & 7 & 7 \\
B, $h=5$ & 1991--1995 & 5 & 4 & 5 & 0 & 4 & 3 \\
B, $h=5$ & 1996--2012 & 17 & 17 & 17 & 0 & 0 & 0 \\
B, $h=5$ & 2013--2019 & 7 & 5 & 7 & 0 & 0 & 0 \\
\bottomrule
\end{tabular}
\end{table}

On Specification A the one-step M1 forecast is at the floor at every origin from 1993 to 2007; at 2008--2012 it forecasts 13--37\,kt from states of 91--169\,kt, and at 2013--2014 it forecasts 105 and 204\,kt against 250 and 299. The mechanism is the constant catch: at the 2000 origin the training-mean catch is 129\,kt (coarse) or 122\,kt (annual) against a state of 34\,kt, and at the 2010 origin 85 or 79\,kt against 91\,kt. On Specification B the training-mean catch is 317\,kt at the 1995 origin and still 201\,kt at the 2020 origin.

The log-RMSE diagnostic of the archived summaries uses $\log\max(\hat S,\epsilon)$ and is dominated by these floor forecasts; over the whole origin sets the floor binds for M1 at 15 of 25 ($h=1$) and 19 of 21 ($h=5$) origins on the Specification A annual pass and at 22 of 59 and 36 of 55 on Specification B, for M1b at 17/25, 19/21, 24/59 and 46/55, and for M2--M4 at between zero and eleven origins per cell. The raw RMSE, not the log column, is the retention score.

\FloatBarrier
\section{Secondary scores}\label{si:secondary}

\begin{table}[!htb]
\centering\small
\caption{Brier score of the hard below-threshold indicator implied by the point forecast ($\hat S < $ threshold against $S <$ threshold), origin-matched. Specification A uses the LRP of 884.6\,kt, below which every origin and target after 1990 lies, so the indicator never varies and persistence scores zero trivially; Specification B uses $\Blim = 276\kt$, which the series crosses in 1977, 1980, 1993 and 2016.}
\label{tab:si_brier}
\begin{tabular}{@{}lrrrr@{}}
\toprule
 & A, $h=1$ & A, $h=5$ & B, $h=1$ & B, $h=5$ \\
\midrule
persistence & 0.000 & 0.000 & 0.068 & 0.291 \\
M1 & 0.040 & 0.048 & 0.051 & 0.455 \\
M1b & 0.000 & 0.048 & 0.153 & 0.473 \\
M2 & 0.080 & 0.143 & 0.068 & 0.364 \\
M3 & 0.080 & 0.095 & 0.017 & 0.309 \\
M4 & 0.080 & 0.143 & 0.068 & 0.382 \\
\bottomrule
\end{tabular}
\end{table}

On Specification B the M3 indicator has the lowest one-step Brier score (0.017 against 0.068 for persistence: one misclassified crossing instead of four), and M1's five-year score (0.455) is the worst because it forecasts the floor through two decades in which the stock sat just below $\Blim$ and then crossed it. The direction-of-change score, the fraction of consecutive targets whose observed and forecast changes agree in sign, is 0.80--1.00 for the structural modules on Specification A and undefined for persistence (its forecast change is zero); it separates modules on an outcome that scarcely varies and carries no retention weight.

An interval-scale comparison is available on Specification B because the assessment publishes 95\% limits for its own predictand, with median relative width 0.44. Writing $D = |\hat S_{t+h|t} - S_{t+h}| / [(\mathrm{hi}_{t+h}-\mathrm{lo}_{t+h})/2]$, the proportion of hindcast errors not exceeding half the published width is 56\% for persistence, 36\% for M1 and 59\% for M3 at $h=1$ ($n=59$), and 18\%, 5\% and 7\% at $h=5$ ($n=55$). These are not coverage probabilities: the published interval expresses uncertainty about true biomass given the assessment's data, re-centring its width on a hindcast gives an interval with no established coverage, and the limits are right-skewed on the biomass scale (67 of 71 years asymmetric by more than 10\%). The statistic indicates only that hindcast discrepancies are of the same order as, or larger than, reported target-state uncertainty, and that M3 reads marginally above persistence on it at one step where the RMSE ordering does not.

\FloatBarrier
\section{Simulation: design and full decomposition}\label{si:sim}

\begin{table}[!htb]
\centering\small
\caption{Data-generating processes. All series have $T=33$ years (the Specification A length) with Gaussian process noise $\sigma\in\{11.8, 33.8\}\kt$ entering the state update; parameters are the archived fixed-window fits. Rolling origins use the eight-year minimum training length of Specification A. M3 and M4 were not simulated as truths.}
\label{tab:si_dgp}
\begin{tabularx}{\textwidth}{@{}lLL@{}}
\toprule
DGP & Generating process & Source of parameters \\
\midrule
D1 & M1 autonomous Schaefer, $r=1.935$, $K=1032.7$, $C\equiv 240$, $S_0 = 900$ & collapse-window coarse fit \\
D2 & M1 autonomous Schaefer, $r=0.458$, $K=500$, $C\equiv 5$, $S_0 = 30$ & recovery-window coarse fit \\
D3 & M2 stock-flow, $r=1.935$, $K=1032.7$, catch 240/120/5 along the Specification A schedule, $S_0=900$ & collapse fit and regime path \\
D4 & M1b depensatory, $r=0.458$, $K=500$, $\sA=15$, $C\equiv 5$, $S_0=30$ & recovery fit with an identifiable threshold \\
D5 & Persistence-true null, $S_{t+1} = S_t + \eta_t$, no production, $C\equiv 0$, $S_0=300$ & -- \\
D6 & Drifting productivity, $r_t = 1.935\,e^{-0.05t} + 0.35$, $K=1032.7$, catch 55\% of the instantaneous $r_tK/4$ & Amendment 1 (outside the class) \\
D7 & Smooth path: noise-free state at $r=0.9$, $K=1032.7$, $C=180$, scored series $=$ state plus Gaussian noise & Amendment 1 (outside the class) \\
\bottomrule
\end{tabularx}
\end{table}

The design sheet (\texttt{SPECIFICATION\_v4.md}) was locked at commit \texttt{d5d9758} on 2026-09-10 and the simulation code was written afterwards to implement it; no design element was changed after execution. Amendment 1 (D6, D7, and a $T=71$ series-length cell for D1 and D5) was locked at commit \texttt{bd0008f} on 2026-09-11; the $T=71$ cell was started and abandoned as infeasible (45.6\,s per rolling pass against 7\,s at $T=33$), so no series-length sensitivity is claimed. A module counts as retained in a replicate if it passes cod Protocol C1--C3, including the 5\% relative-margin rule, at both horizons exactly as in the empirical rule. The completed v68 rerun covers 2,800 seeded DGP/noise series (14 cells $\times$ 200 replicates), scored against five modules, yielding 14,000 raw module rows. Replicate IDs 0--199 were processed in four disjoint 50-ID batches; seeds use the first eight bytes of SHA-256 over a fixed namespace, DGP cell, noise level to one decimal place, and replicate ID, with NumPy's PCG64 generator. All four batch manifests record source and output hashes; no replicate failed or produced a non-finite score. The estimates remain conditional on these fixed DGPs and parameter settings.

\begin{table}[!htb]
\centering\small
\caption{Cod Protocol C operating characteristics. Panel A decomposes power loss averaged over the two process-noise levels: the persistence-only hurdle applies the 5\% band at both horizons but omits the next-simpler comparator. Panel B gives each cell's observed rate and 95\% Wilson interval; D1--D4 are truth-matching-module retention, while D5--D7 are any-module retention. Wilson intervals quantify Monte Carlo uncertainty only, not uncertainty in the fitted models or data-generating processes.}
\label{tab:si_power}
\noindent\textit{Panel A. Persistence-only versus full-rule power.}\par
\scriptsize\setlength{\tabcolsep}{2.5pt}
\begin{tabular}{@{}llrrrr@{}}
\toprule
DGP & truth & persistence-only hurdle & full C rule & removed by comparator & conditional on hurdle \\
\midrule
D1 & M1 & 0.9625 & 0.9625 & 0.0000 & 0\% \\
D2 & M1 & 0.3850 & 0.3850 & 0.0000 & 0\% \\
D3 & M2 & 0.3300 & 0.1175 & 0.2125 & 64\% \\
D4 & M1b & 0.1625 & 0.0075 & 0.1550 & 95\% \\
\bottomrule
\end{tabular}

\par\medskip
\noindent\textit{Panel B. Cell-specific retention rate (count/200; rate [95\% Wilson interval]).}\par
\scriptsize
\begin{tabularx}{\textwidth}{@{}l>{\centering\arraybackslash}X>{\centering\arraybackslash}X@{}}
\toprule
DGP & $\sigma=11.8\,\kt$ & $\sigma=33.8\,\kt$ \\
\midrule
D1 (M1 true) & 195/200; 0.975 [0.943, 0.989] & 190/200; 0.950 [0.910, 0.973] \\
D2 (M1 true) & 132/200; 0.660 [0.592, 0.722] & 22/200; 0.110 [0.074, 0.161] \\
D3 (M2 true) & 21/200; 0.105 [0.070, 0.155] & 26/200; 0.130 [0.090, 0.184] \\
D4 (M1b true) & 0/200; 0.000 [0.000, 0.019] & 3/200; 0.015 [0.005, 0.043] \\
D5 (persistence null) & 5/200; 0.025 [0.011, 0.057] & 15/200; 0.075 [0.046, 0.120] \\
D6 (drifting productivity) & 144/200; 0.720 [0.654, 0.778] & 154/200; 0.770 [0.707, 0.823] \\
D7 (smooth path + observation error) & 194/200; 0.970 [0.936, 0.986] & 190/200; 0.950 [0.910, 0.973] \\
\bottomrule
\end{tabularx}
\end{table}

\begin{table}[!htb]
\centering\small
\caption{Descriptive rank-baseline check: proportion of 400 replicates (200 at each noise level) in which the generating module has the lowest one-step RMSE among the five modules; the nominal uniform-rank baseline is 0.20. Labels are descriptive comparisons, not formal tests of a random-ranking null.}
\label{tab:si_chance}
\begin{tabular}{@{}lrl@{}}
\toprule
DGP & true module ranks first & against chance \\
\midrule
D1 (M1 true) & 0.623 & far above 0.20 \\
D2 (M1 true) & 0.740 & far above 0.20 \\
D3 (M2 true) & 0.028 & below 0.20 \\
D4 (M1b true) & 0.245 & slightly above 0.20 \\
\bottomrule
\end{tabular}
\end{table}

Panel A separates the persistence hurdle from comparator attrition. Under D3, the persistence-only hurdle is passed in 12.0\% of low-noise and 54.0\% of high-noise series; the full rule retains true M2 in 10.5\% and 13.0\%. The comparator removes 3 of 24 and 82 of 108 persistence-hurdle survivors (12.5\% and 75.9\%), respectively. Under D4 the persistence-only rates are 28.5\% and 4.0\%, but the full rule retains M1b in 0.0\% and 1.5\%; the comparator removes 57/57 and 5/8 survivors. Any module is retained under D3 in 105/200 and 122/200 series (95\% Wilson intervals [0.456, 0.593] and [0.541, 0.675]), and under D4 in 61/200 and 11/200 ([0.245, 0.372] and [0.031, 0.096]), illustrating that low truth-specific power is not equivalent to no model being selected. Across D3's 400 simulated series, mean one-step RMSE is 85.5\,kt for the true M2, versus 77.5\,kt for M3 and 80.3\,kt for M1b (78.2\,kt for M4); estimation noise can therefore make the correct structure score worse than its siblings. Across D1--D4, non-truth modules are retained in 287 of 6,400 non-truth module--replicate pairs (4.5\%), or 0.179 false module retentions per simulated series. Under the persistence-true null, 20 of 2,000 module--replicate pairs are retained (1.0\%), averaging 0.050 false modules per series. Under D6, M1 is retained in 144/200 and 154/200 replicates at the two noise levels (298 M1 retentions among 300 module-retention events); under D7 the counts are 194/200 and 189/200 (383 of 394 events). A constant-productivity map thus out-predicts persistence on smoothly drifting or smoothly observed series often enough to clear both hurdles. Table~\ref{tab:si_rules} keeps the earlier alternative-rule summaries as reported, but its Protocol C reference row is recomputed from the completed v68 archive. The calculation code or complete formula for the alternative rules was not located for this edition; those alternatives are therefore not independently reproduced and do not inform the primary-rule conclusions.

The earlier archived analysis reported two candidate diagnostics for predicting from the data alone whether the rule has power: the dispersion of scores across modules (median coefficient of variation) correlated with true-module power at Spearman $\rho=0.52$, while the best-module margin over persistence correlated at $\rho=0.88$ (Pearson $r=0.94$). These summaries are retained as historical reports, not independently reproduced for v68: their calculation code and full cell-level formula were not located. The second diagnostic would in any case reuse the same out-of-sample scores consumed by the rule and cannot serve as an independent gate. Identifying a pre-score indicator of power remains open.

\FloatBarrier

\begin{table}[!htb]
\centering\small
\caption{Alternative decision-rule summaries. The Protocol C reference row uses the full v68 rerun; the other rules are carried forward from the earlier archived analysis and were not independently reproduced because their calculation code or complete formula was not located. Mean power is averaged over D1--D4 and both noise levels; specificity is one minus the any-module false-retention rate under D5. The earlier D6--D7 comparison is likewise historical and not recalculated here.}
\label{tab:si_rules}
\begin{tabular}{@{}lrr@{}}
\toprule
Decision rule & mean power & specificity \\
\midrule
Protocol C as adopted (C1--C3, 5\% band) & 0.368 & 0.950 \\
Without the comparator hurdle (C2, C3 only) & 0.476 & 0.972 \\
Persistence only, $h=1$, 5\% band & 0.562 & 0.955 \\
Persistence only, any margin & 0.542 & 0.765 \\
Full rule, no tie band & 0.436 & 0.772 \\
Full rule, 10\% band & 0.337 & 0.998 \\
MASE $<1$ against the naive benchmark & 0.651 & 0.675 \\
Information criterion, $n\log\mathrm{MSE} + 2k$ & 0.509 & 0.992 \\
\bottomrule
\end{tabular}
\end{table}

\section{Decomposition of the lagged-start deficit}\label{si:m4}

M4 differs from M3 in two ways: it starts the path from $S_{t-1}$, and it therefore integrates one more catch year. A delayed-persistence control, $\hat S_{t+h} = S_{t-1}$, isolates the information delay. Its one-step RMSE is 184.4\,kt on Specification A and 158.0\,kt on Specification B, against 98.0 and 84.4 for ordinary persistence; the delay alone therefore costs 86.4 of M4's 97.5-kt deficit against persistence on A and 73.6 of 121.3 on B. The lagged start is an information-delay penalty, not a dynamical one, and M4 is not retained for the same reason that a one-year-stale persistence forecast would not be.

\FloatBarrier
\section{The outcome years 2016--2021}\label{si:outcome}

Specification A ends in 2015. The continuation of the same assessment family is published in the RAM Legacy Stock Assessment Database (v4.66, stock COD2J3KL), and xteNCAM reads the same years independently (Table~\ref{tab:si_outcome}). Every year 2016--2021 lies above its 2015 origin in both vintages, so the window contains no event of the kind the collapse test probes, and persistence from the 2015 origin errs low throughout: MAE 129.2\,kt and RMSE 133.3\,kt on the RAM vintage (107.5\,kt MAE from the Table A2 origin of 298.65\,kt), the same order as its scored one-step RMSE. The two vintages agree on the outcome years to within 1--26\,kt (mean 16.8), while on the historical overlap they do not (S11). The restatement of the Edwards J-17 scores under Protocol C, which the single-system edition carried here, is now in main-text Section 4.3.

\begin{table}[!htb]
\centering\small
\caption{Outcome years on the Specification A horizon. Persistence error is $|S_t - S_{2015}|$ from the RAM-vintage 2015 origin of 277\,kt; the xteNCAM column is the revised model's reading of the same years.}
\label{tab:si_outcome}
\begin{tabular}{@{}lrrrr@{}}
\toprule
Year & RAM v4.66 (kt) & xteNCAM (kt) & $|$difference$|$ & persistence error (kt) \\
\midrule
2016 & 340 & 339 & 1 & 63 \\
2017 & 433 & 451 & 18 & 156 \\
2018 & 394 & 369 & 25 & 117 \\
2019 & 419 & 400 & 19 & 142 \\
2020 & 440 & 414 & 26 & 163 \\
2021 & 411 & 423 & 12 & 134 \\
\bottomrule
\end{tabular}
\end{table}
